\section{Related Work}



ID estimation methods can be divided into two broad categories: global and local \citep{camastra2016intrinsic}. Global methods aim to give a single estimate of the dimensionality of the entire dataset, which however discards the nuanced manifold structure when the data lies on a union of different manifolds (which is often the case for real-world datasets). 

On the contrary, the local methods \citep{carter2009local,kleindessner2015dimensionality,levina2004maximum,hino2017local,camastra2016intrinsic, rozza2012novel,CERUTI20142569,camastra2002} try to estimate the local ID of the data manifold at an arbitrary point.
This approach gives more insight into the nature of the dataset, and provides more options to summarize the dimensionality of the manifold than the global perspective.
A detailed overview of the methods used for global and local ID estimation is provided by \citet{camastra2016intrinsic}, and for a good review of the local ID estimation methods we refer the reader to \citet{johnsson2014low}. 
We list all the algorithms we compare to in Table~\ref{tab:skdim} in the appendices.







